% GENERATED FILE. Edit canonical lessons, then run npm run generate:books.
% canonical-source-hash: fnv1a64:0e9c74e5f9431831
% canonical-lessons: MR-C350-svikarne, MR-C350-manya-karne, MR-C350-grahya-dharne, MR-C350-dos-dene, MR-C350-arop-karne

\chapter{Accepting and Blaming}
\label{ch:mr-accepting-and-blaming}
\hlchaptermodality{350}

\begin{chapteropening}
\textbf{By the end of this chapter:} \emph{I can accept an argument, admit a fault or lay blame.}

Complete the last lesson of chapter 350: I can accept an argument, admit a fault or lay blame.
\end{chapteropening}

\section[\texorpdfstring{\mr{स्वीकारणे} (svīkārṇe)}{स्वीकारणे (svīkārṇe)}]{\mr{स्वीकारणे} (svīkārṇe) — to accept}

\label{lesson:MR-C350-svikarne}



\begin{quote}
Before the new one: say the Marathi for an alternative, then the Marathi for to unravel.
\end{quote}

\subsection*{You'll want to know: \mr{स्वीकारणे}}
\textbf{\mr{स्वीकारणे}} — \emph{svīkārṇe} — \textquotedblleft{}to accept\textquotedblright{}.

\begin{quote}
\mr{त्याने} \mr{नवीन} \mr{नोकरी} \mr{आनंदाने} \mr{स्वीकारली}.
\end{quote}

\begin{quote}
\emph{tyāne navīn nokrī ānandāne svīkārlī.}
\end{quote}

\begin{quote}
He gladly accepted the new job.
\end{quote}

Past, with -\mr{ने} (he …): \mr{त्याने} \mr{स्वीकारलं}, \emph{tyāne svīkārlaṁ}.

Transitive: the past agrees with what is accepted, \mr{नोकरी} \mr{स्वीकारली}.

\begin{cousinweb}
\textbf{\mr{स्वीकारणे}} (\emph{svīkārṇe}) is built on the tatsama \mr{स्वीकार}, Sanskrit \emph{svīkāra}, \textquotedblleft{}making one's own, accepting\textquotedblright{} (Monier-Williams), with the verb ending -\mr{णे}.
\end{cousinweb}

\begin{grammarlens}[title={What you've built}]
One more word for giving reasons and opinions.
\end{grammarlens}

\subsection*{Your turn}
\begin{itemize}
\raggedright
  \item \emph{Say it:} \emph{svīkārṇe}
  \item \emph{Say it:} \emph{svīkārṇe}, once more
  \item \emph{Say it:} \emph{tyāne navīn nokrī ānandāne svīkārlī.} — He gladly accepted the new job.
  \item \emph{Recall:} say the Marathi for an alternative, then the Marathi for to unravel, then say \emph{svīkārṇe} again
\end{itemize}

\begin{tcolorbox}[breakable,colback=teal!4,colframe=teal!35!black,title={Before you move on}]
What does \mr{स्वीकारणे} mean? (\textquotedblleft{}To accept\textquotedblright{}.) Say it once more.
\end{tcolorbox}
\section[\texorpdfstring{\mr{मान्य} \mr{करणे} (mānya karṇe)}{मान्य करणे (mānya karṇe)}]{\mr{मान्य} \mr{करणे} (mānya karṇe) — to acknowledge; to agree to}

\label{lesson:MR-C350-manya-karne}



\begin{quote}
Before the new one: say the Marathi for to unravel, then the Marathi for to accept.
\end{quote}

\subsection*{You'll want to know: \mr{मान्य} \mr{करणे}}
\textbf{\mr{मान्य} \mr{करणे}} — \emph{mānya karṇe} — \textquotedblleft{}to acknowledge; to agree to\textquotedblright{}.

\begin{quote}
\mr{शेवटी} \mr{त्याने} \mr{आपली} \mr{चूक} \mr{मान्य} \mr{केली}.
\end{quote}

\begin{quote}
\emph{śevṭī tyāne āplī cūk mānya kelī.}
\end{quote}

\begin{quote}
In the end he acknowledged his mistake.
\end{quote}

Past, with -\mr{ने} (he …): \mr{त्याने} \mr{मान्य} \mr{केलं}, \emph{tyāne mānya kelaṁ}.

\mr{मान्य} \mr{असणे}, to be acceptable: \mr{मला} \mr{हे} \mr{मान्य} \mr{आहे}.

\begin{cousinweb}
\textbf{\mr{मान्य} \mr{करणे}} (\emph{mānya karṇe}) joins the tatsama \mr{मान्य}, Sanskrit \emph{mānya}, \textquotedblleft{}to be respected or honoured\textquotedblright{} (Monier-Williams), to \mr{करणे}, \textquotedblleft{}to make\textquotedblright{}.
\end{cousinweb}

\begin{grammarlens}[title={What you've built}]
One more word for giving reasons and opinions.
\end{grammarlens}

\subsection*{Your turn}
\begin{itemize}
\raggedright
  \item \emph{Say it:} \emph{mānya karṇe}
  \item \emph{Say it:} \emph{mānya karṇe}, once more
  \item \emph{Say it:} \emph{śevṭī tyāne āplī cūk mānya kelī.} — In the end he acknowledged his mistake.
  \item \emph{Recall:} say the Marathi for to unravel, then the Marathi for to accept, then say \emph{mānya karṇe} again
\end{itemize}

\begin{tcolorbox}[breakable,colback=teal!4,colframe=teal!35!black,title={Before you move on}]
What does \mr{मान्य} \mr{करणे} mean? (\textquotedblleft{}To acknowledge; to agree to\textquotedblright{}.) Say it once more.
\end{tcolorbox}
\section[\texorpdfstring{\mr{ग्राह्य} \mr{धरणे} (grāhya dharṇe)}{ग्राह्य धरणे (grāhya dharṇe)}]{\mr{ग्राह्य} \mr{धरणे} (grāhya dharṇe) — to treat as valid}

\label{lesson:MR-C350-grahya-dharne}



\begin{quote}
Before the new one: say the Marathi for to accept, then the Marathi for to acknowledge.
\end{quote}

\subsection*{You'll want to know: \mr{ग्राह्य} \mr{धरणे}}
\textbf{\mr{ग्राह्य} \mr{धरणे}} — \emph{grāhya dharṇe} — \textquotedblleft{}to treat as valid\textquotedblright{}.

\begin{quote}
\mr{शिक्षकांनी} \mr{त्याचं} \mr{दुसरं} \mr{उत्तरही} \mr{ग्राह्य} \mr{धरलं}.
\end{quote}

\begin{quote}
\emph{śikṣakānnī tyācaṁ dusraṁ uttarhī grāhya dharlaṁ.}
\end{quote}

\begin{quote}
The teacher accepted his second answer as valid too.
\end{quote}

Past, with -\mr{ने} (he …): \mr{त्याने} \mr{ग्राह्य} \mr{धरलं}, \emph{tyāne grāhya dharlaṁ}.

\mr{धरणे}, to hold: to hold something as acceptable.

\begin{cousinweb}
\textbf{\mr{ग्राह्य} \mr{धरणे}} (\emph{grāhya dharṇe}) joins the tatsama \mr{ग्राह्य}, Sanskrit \emph{grāhya}, \textquotedblleft{}to be seized, to be accepted\textquotedblright{} (Monier-Williams), to \mr{धरणे}, \textquotedblleft{}to hold\textquotedblright{}.
\end{cousinweb}

\begin{grammarlens}[title={What you've built}]
One more word for giving reasons and opinions.
\end{grammarlens}

\subsection*{Your turn}
\begin{itemize}
\raggedright
  \item \emph{Say it:} \emph{grāhya dharṇe}
  \item \emph{Say it:} \emph{grāhya dharṇe}, once more
  \item \emph{Say it:} \emph{śikṣakānnī tyācaṁ dusraṁ uttarhī grāhya dharlaṁ.} — The teacher accepted his second answer as valid too.
  \item \emph{Recall:} say the Marathi for to accept, then the Marathi for to acknowledge, then say \emph{grāhya dharṇe} again
\end{itemize}

\begin{tcolorbox}[breakable,colback=teal!4,colframe=teal!35!black,title={Before you move on}]
What does \mr{ग्राह्य} \mr{धरणे} mean? (\textquotedblleft{}To treat as valid\textquotedblright{}.) Say it once more.
\end{tcolorbox}
\section[\texorpdfstring{\mr{दोष} \mr{देणे} (doṣ deṇe)}{दोष देणे (doṣ deṇe)}]{\mr{दोष} \mr{देणे} (doṣ deṇe) — to blame}

\label{lesson:MR-C350-dos-dene}



\begin{quote}
Before the new one: say the Marathi for to acknowledge, then the Marathi for to treat as valid.
\end{quote}

\subsection*{You'll want to know: \mr{दोष} \mr{देणे}}
\textbf{\mr{दोष} \mr{देणे}} — \emph{doṣ deṇe} — \textquotedblleft{}to blame\textquotedblright{}.

\begin{quote}
\mr{तो} \mr{आपल्या} \mr{अपयशासाठी} \mr{नशिबाला} \mr{दोष} \mr{देतो}.
\end{quote}

\begin{quote}
\emph{to āpalyā apayśāsāṭhī naśibālā doṣ deto.}
\end{quote}

\begin{quote}
He blames fate for his failure.
\end{quote}

Past, with -\mr{ने} (he …): \mr{त्याने} \mr{दोष} \mr{दिला}, \emph{tyāne doṣ dilā}.

The one blamed takes -\mr{ला}; the fault takes -\mr{साठी} or -\mr{बद्दल}.

\begin{cousinweb}
\textbf{\mr{दोष} \mr{देणे}} (\emph{doṣ deṇe}) joins the tatsama \mr{दोष}, Sanskrit \emph{doṣa}, \textquotedblleft{}fault, defect\textquotedblright{} (Monier-Williams), to \mr{देणे}, \textquotedblleft{}to give\textquotedblright{}.
\end{cousinweb}

\begin{grammarlens}[title={What you've built}]
One more word for giving reasons and opinions.
\end{grammarlens}

\subsection*{Your turn}
\begin{itemize}
\raggedright
  \item \emph{Say it:} \emph{doṣ deṇe}
  \item \emph{Say it:} \emph{doṣ deṇe}, once more
  \item \emph{Say it:} \emph{to āpalyā apayśāsāṭhī naśibālā doṣ deto.} — He blames fate for his failure.
  \item \emph{Recall:} say the Marathi for to acknowledge, then the Marathi for to treat as valid, then say \emph{doṣ deṇe} again
\end{itemize}

\begin{tcolorbox}[breakable,colback=teal!4,colframe=teal!35!black,title={Before you move on}]
What does \mr{दोष} \mr{देणे} mean? (\textquotedblleft{}To blame\textquotedblright{}.) Say it once more.
\end{tcolorbox}
\section[\texorpdfstring{\mr{आरोप} \mr{करणे} (ārop karṇe)}{आरोप करणे (ārop karṇe)}]{\mr{आरोप} \mr{करणे} (ārop karṇe) — to accuse; to allege}

\label{lesson:MR-C350-arop-karne}



\begin{quote}
Before the new one: say the Marathi for to treat as valid, then the Marathi for to blame.
\end{quote}

\subsection*{You'll want to know: \mr{आरोप} \mr{करणे}}
\textbf{\mr{आरोप} \mr{करणे}} — \emph{ārop karṇe} — \textquotedblleft{}to accuse; to allege\textquotedblright{}.

\begin{quote}
\mr{शेजारी} \mr{त्याच्यावर} \mr{चोरीचा} \mr{आरोप} \mr{करत} \mr{होता}.
\end{quote}

\begin{quote}
\emph{śejārī tyācyāvar corīcā ārop karat hotā.}
\end{quote}

\begin{quote}
The neighbour was accusing him of theft.
\end{quote}

Past, with -\mr{ने} (he …): \mr{त्याने} \mr{आरोप} \mr{केला}, \emph{tyāne ārop kelā}.

The accused takes -\mr{वर} and the charge -\mr{चा}: \mr{त्याच्यावर} \mr{चोरीचा} \mr{आरोप}.

\begin{cousinweb}
\textbf{\mr{आरोप} \mr{करणे}} (\emph{ārop karṇe}) joins the tatsama \mr{आरोप}, Sanskrit \emph{āropa}, \textquotedblleft{}placing upon, imputing\textquotedblright{} (Monier-Williams), to \mr{करणे}, \textquotedblleft{}to do\textquotedblright{}.
\end{cousinweb}

\begin{grammarlens}[title={What you've built}]
One more word for giving reasons and opinions.
\end{grammarlens}

\subsection*{Your turn}
\begin{itemize}
\raggedright
  \item \emph{Say it:} \emph{ārop karṇe}
  \item \emph{Say it:} \emph{ārop karṇe}, once more
  \item \emph{Say it:} \emph{śejārī tyācyāvar corīcā ārop karat hotā.} — The neighbour was accusing him of theft.
  \item \emph{Recall:} say the Marathi for to treat as valid, then the Marathi for to blame, then say \emph{ārop karṇe} again
\end{itemize}

\begin{tcolorbox}[breakable,colback=teal!4,colframe=teal!35!black,title={Before you move on}]
What does \mr{आरोप} \mr{करणे} mean? (\textquotedblleft{}To accuse; to allege\textquotedblright{}.) Say it once more.
\end{tcolorbox}
